% 已确认的经典蓝色 Madrid 主题
\usetheme{Madrid}
\usefonttheme{professionalfonts}
\definecolor{AcademicBlue}{RGB}{51,51,178}
\definecolor{BlueDark}{RGB}{26,26,89}
\definecolor{BlueMid}{RGB}{38,38,134}
\definecolor{NoteGray}{RGB}{105,105,105}
\definecolor{GridGray}{RGB}{220,220,230}
\setbeamercolor{structure}{fg=AcademicBlue}
\setbeamercolor{palette primary}{fg=white,bg=AcademicBlue}
\setbeamercolor{palette secondary}{fg=white,bg=BlueMid}
\setbeamercolor{palette tertiary}{fg=white,bg=BlueDark}
\setbeamercolor{frametitle}{fg=white,bg=AcademicBlue}
\setbeamercolor{title}{fg=white,bg=AcademicBlue}
\setbeamercolor{subtitle}{fg=white}
\setbeamercolor{normal text}{fg=black,bg=white}
\setbeamercolor{block title}{fg=white,bg=AcademicBlue}
\setbeamercolor{block body}{fg=black,bg=AcademicBlue!7!white}
\setbeamercolor{author in head/foot}{fg=white,bg=BlueDark}
\setbeamercolor{title in head/foot}{fg=white,bg=BlueMid}
\setbeamercolor{date in head/foot}{fg=white,bg=AcademicBlue}
\setbeamerfont{title}{size=\Large,series=\mdseries}
\setbeamerfont{subtitle}{size=\normalsize,series=\mdseries}
\setbeamerfont{frametitle}{size=\Large,series=\mdseries}
\setbeamerfont{block title}{size=\normalsize,series=\mdseries}
\setbeamerfont{author}{size=\normalsize}
\setbeamerfont{institute}{size=\small}
\setbeamerfont{date}{size=\normalsize}
\setbeamersize{text margin left=5mm,text margin right=5mm}
\setbeamertemplate{navigation symbols}{}
\setbeamertemplate{headline}{}
\setbeamertemplate{enumerate items}[ball]
\setbeamertemplate{items}[ball]
\setbeamertemplate{blocks}[rounded][shadow=true]
% 部分旧版PGF的dvipdfmx驱动在保存阴影对象编号时使用了错误的let展开。
% 下方仅对具有该对象计数器的dvipdfmx类后端应用等价的资源登记定义。
\makeatletter
\AtBeginDocument{%
  \@ifundefined{pgfsys@objnum}{}{%
    \def\pgf@sys@pdf@install@mask#1{%
      \pgfutil@addpdfresource@extgs{%
        /pgfsmask\csname pgfsmaskxform@#1\endcsname\space <<
        /SMask << /S /Luminosity /G @pgfsmask\csname pgfsmaskxform@#1\endcsname >> >>}%
      \expandafter\xdef\csname pgfsmask@#1\endcsname{\csname pgfsmaskxform@#1\endcsname}%
    }%
  }%
}
\makeatother
\setlength{\parindent}{0pt}
\setbeamertemplate{footline}{%
  \leavevmode\hbox{%
  \begin{beamercolorbox}[wd=.25\paperwidth,ht=2.1ex,dp=.85ex,center]{author in head/foot}%
    \fontsize{5}{6}\selectfont\bfseries\insertshortauthor\quad 空间转录组学入门
  \end{beamercolorbox}%
  \begin{beamercolorbox}[wd=.416666\paperwidth,ht=2.1ex,dp=.85ex,center]{title in head/foot}%
    \fontsize{5}{6}\selectfont\bfseries\insertsubsectionhead
  \end{beamercolorbox}%
  \begin{beamercolorbox}[wd=.333334\paperwidth,ht=2.1ex,dp=.85ex]{date in head/foot}%
    \fontsize{5}{6}\selectfont\bfseries\hspace{5mm}\insertshortdate
    \hfill\insertframenumber\,/\,\inserttotalframenumber\hspace{1.5mm}
  \end{beamercolorbox}}}
\hypersetup{colorlinks=true,urlcolor=NoteGray,linkcolor=white,citecolor=NoteGray}
\newcommand{\source}[1]{\vfill\par{\fontsize{6}{8}\selectfont\color{NoteGray}\raggedright #1\par}}
\newcommand{\teaching}{\source{教学构造：沿用原讲义示例，不是真实实验数据。}}
\newcommand{\bluelabel}[1]{{\color{AcademicBlue}#1}\par\smallskip}
\newcommand{\keypoint}[2]{\begin{block}{#1}#2\end{block}}
\newcommand{\wideitems}{\setlength{\itemsep}{.75em}}
\newcommand{\compactitems}{\setlength{\itemsep}{.35em}}
\newcolumntype{Y}{>{\raggedright\arraybackslash}X}
\newcommand{\codetext}[1]{\texttt{#1}}
